\section{fill}

Kreyszig 9.9 Spectral Representation of Bounded Self-Adjoint Linear Operators
1 First some easy examples
2
3
4 Let us apply the result to the finite-dimensional case and see in what we recover the og spectral theorem
5
6 In what way do we get the compact spectral theorem?
7 This is the position operator from quantum mechanics? What would happen in general for y(t) = v(t) x(t) where v is continuous
8 What would happen in general for a map (xi1, xi2, ... ) -> (a1 xi1, a2 xi2, ...)   where a1, a2,... are real and bounded (necessary to make T self-adjoint and bounded)
9
10

\begin{exercise}{1}
fill
\end{exercise}
\begin{proof}
fill
\end{proof} 

\begin{exercise}{2}
fill
\end{exercise}
\begin{proof}
fill
\end{proof} 

\begin{exercise}{3}
fill
\end{exercise}
\begin{proof}
fill
\end{proof} 

\begin{exercise}{4}
fill
\end{exercise}
\begin{proof}
fill
\end{proof} 

\begin{exercise}{5}
fill
\end{exercise}
\begin{proof}
fill
\end{proof} 

\begin{exercise}{6}
fill
\end{exercise}
\begin{proof}
fill
\end{proof} 

\begin{exercise}{7}
fill
\end{exercise}
\begin{proof}
fill
\end{proof} 

\begin{exercise}{8}
fill
\end{exercise}
\begin{proof}
fill
\end{proof} 

\begin{exercise}{9}
fill
\end{exercise}
\begin{proof}
fill
\end{proof} 

\begin{exercise}{10}
fill
\end{exercise}
\begin{proof}
fill
\end{proof} 
